% GENERATED FILE. Edit canonical lessons, then run npm run generate:books.
% canonical-source-hash: fnv1a64:4ef9d0744865e3c2
% canonical-lessons: TA-C511-errumati, TA-C511-irakkumati, TA-C511-vanikam, TA-C511-tolilatipar, TA-C511-kuttali

\chapter{Imports and Exports}
\label{ch:ta-imports-and-exports}
\hlchaptermodality{511}

\begin{chapteropening}
\textbf{By the end of this chapter:} \emph{I can talk about trade between countries.}

Complete the last lesson of chapter 511: I can talk about trade between countries.
\end{chapteropening}

\section[\texorpdfstring{\ta{ஏற்றுமதி} (ēṟṟumati)}{ஏற்றுமதி (ēṟṟumati)}]{\ta{ஏற்றுமதி} (ēṟṟumati) — exports}

\label{lesson:TA-C511-errumati}



\begin{quote}
Before the new one: say the Tamil for to stack, then the Tamil for to take measurements.
\end{quote}

\subsection*{You'll want to know: \ta{ஏற்றுமதி}}
\textbf{\ta{ஏற்றுமதி}} — \emph{ēṟṟumati} — \textquotedblleft{}exports\textquotedblright{}.

\begin{quote}
\ta{தேயிலை} \ta{இந்த} \ta{மாவட்டத்தின்} \ta{முக்கியமான} \ta{ஏற்றுமதி}.
\end{quote}

\begin{quote}
\emph{tēyilai inta māvaṭṭattiṉ mukkiyamāṉa ēṟṟumati.}
\end{quote}

\begin{quote}
Tea is the main export of this district.
\end{quote}

With \ta{செய்} (\emph{cey}) it becomes a verb: \ta{ஏற்றுமதி} \ta{செய்}, to export.

\begin{cousinweb}
\textbf{\ta{ஏற்றுமதி}} (\emph{ēṟṟumati}) is built on the native \ta{ஏற்று} (\emph{ēṟṟu}), \textquotedblleft{}to load, to put on board\textquotedblright{}; the origin of the second part, -\ta{மதி}, is not settled.
\end{cousinweb}

\begin{grammarlens}[title={What you've built}]
One more word for describing what you lived and felt.
\end{grammarlens}

\subsection*{Your turn}
\begin{itemize}
\raggedright
  \item \emph{Say it:} \emph{ēṟṟumati}
  \item \emph{Say it:} \emph{ēṟṟumati}, once more
  \item \emph{Say it:} \emph{tēyilai inta māvaṭṭattiṉ mukkiyamāṉa ēṟṟumati.} — Tea is the main export of this district.
  \item \emph{Recall:} say the Tamil for to stack, then the Tamil for to take measurements, then say \emph{ēṟṟumati} again
\end{itemize}

\begin{tcolorbox}[breakable,colback=teal!4,colframe=teal!35!black,title={Before you move on}]
What does \ta{ஏற்றுமதி} mean? (\textquotedblleft{}Exports\textquotedblright{}.) Say it once more.
\end{tcolorbox}
\section[\texorpdfstring{\ta{இறக்குமதி} (iṟakkumati)}{இறக்குமதி (iṟakkumati)}]{\ta{இறக்குமதி} (iṟakkumati) — imports}

\label{lesson:TA-C511-irakkumati}



\begin{quote}
Before the new one: say the Tamil for to take measurements, then the Tamil for exports.
\end{quote}

\subsection*{You'll want to know: \ta{இறக்குமதி}}
\textbf{\ta{இறக்குமதி}} — \emph{iṟakkumati} — \textquotedblleft{}imports\textquotedblright{}.

\begin{quote}
\ta{நம்} \ta{நாடு} \ta{எண்ணெயை} \ta{இறக்குமதி} \ta{செய்கிறது}.
\end{quote}

\begin{quote}
\emph{nam nāṭu eṇṇeyai iṟakkumati ceykiṟatu.}
\end{quote}

\begin{quote}
Our country imports oil.
\end{quote}

The pair of \ta{ஏற்றுமதி} (\emph{ēṟṟumati}): goods unloaded, brought in. With \ta{செய்} it is a verb.

\begin{cousinweb}
\textbf{\ta{இறக்குமதி}} (\emph{iṟakkumati}) is built on the native \ta{இறக்கு} (\emph{iṟakku}), \textquotedblleft{}to unload, to bring down\textquotedblright{}; as with \ta{ஏற்றுமதி}, the origin of -\ta{மதி} is not settled.
\end{cousinweb}

\begin{grammarlens}[title={What you've built}]
One more word for describing what you lived and felt.
\end{grammarlens}

\subsection*{Your turn}
\begin{itemize}
\raggedright
  \item \emph{Say it:} \emph{iṟakkumati}
  \item \emph{Say it:} \emph{iṟakkumati}, once more
  \item \emph{Say it:} \emph{nam nāṭu eṇṇeyai iṟakkumati ceykiṟatu.} — Our country imports oil.
  \item \emph{Recall:} say the Tamil for to take measurements, then the Tamil for exports, then say \emph{iṟakkumati} again
\end{itemize}

\begin{tcolorbox}[breakable,colback=teal!4,colframe=teal!35!black,title={Before you move on}]
What does \ta{இறக்குமதி} mean? (\textquotedblleft{}Imports\textquotedblright{}.) Say it once more.
\end{tcolorbox}
\section[\texorpdfstring{\ta{வணிகம்} (vaṇikam)}{வணிகம் (vaṇikam)}]{\ta{வணிகம்} (vaṇikam) — business; commerce}

\label{lesson:TA-C511-vanikam}



\begin{quote}
Before the new one: say the Tamil for exports, then the Tamil for imports.
\end{quote}

\subsection*{You'll want to know: \ta{வணிகம்}}
\textbf{\ta{வணிகம்}} — \emph{vaṇikam} — \textquotedblleft{}business; commerce\textquotedblright{}.

\begin{quote}
\ta{என்} \ta{மாமா} \ta{துணி} \ta{வணிகம்} \ta{செய்கிறார்}.
\end{quote}

\begin{quote}
\emph{eṉ māmā tuṇi vaṇikam ceykiṟār.}
\end{quote}

\begin{quote}
My uncle is in the cloth business.
\end{quote}

A formal word; in speech people often say \ta{வியாபாரம்} (\emph{viyāpāram}).

\begin{cousinweb}
\textbf{\ta{வணிகம்}} (\emph{vaṇikam}) is borrowed from Sanskrit \emph{vaṇij}, \textquotedblleft{}a trader\textquotedblright{}, also \textquotedblleft{}trade\textquotedblright{} (CDIAL 11230), with the Tamil noun ending -\ta{அம்}.
\end{cousinweb}

\begin{grammarlens}[title={What you've built}]
One more word for describing what you lived and felt.
\end{grammarlens}

\subsection*{Your turn}
\begin{itemize}
\raggedright
  \item \emph{Say it:} \emph{vaṇikam}
  \item \emph{Say it:} \emph{vaṇikam}, once more
  \item \emph{Say it:} \emph{eṉ māmā tuṇi vaṇikam ceykiṟār.} — My uncle is in the cloth business.
  \item \emph{Recall:} say the Tamil for exports, then the Tamil for imports, then say \emph{vaṇikam} again
\end{itemize}

\begin{tcolorbox}[breakable,colback=teal!4,colframe=teal!35!black,title={Before you move on}]
What does \ta{வணிகம்} mean? (\textquotedblleft{}Business; commerce\textquotedblright{}.) Say it once more.
\end{tcolorbox}
\section[\texorpdfstring{\ta{தொழிலதிபர்} (toḻilatipar)}{தொழிலதிபர் (toḻilatipar)}]{\ta{தொழிலதிபர்} (toḻilatipar) — a businessman; an industrialist}

\label{lesson:TA-C511-tolilatipar}



\begin{quote}
Before the new one: say the Tamil for imports, then the Tamil for business.
\end{quote}

\subsection*{You'll want to know: \ta{தொழிலதிபர்}}
\textbf{\ta{தொழிலதிபர்}} — \emph{toḻilatipar} — \textquotedblleft{}a businessman; an industrialist\textquotedblright{}.

\begin{quote}
\ta{அந்தத்} \ta{தொழிலதிபர்} \ta{ஒரு} \ta{புதிய} \ta{தொழிற்சாலையைத்} \ta{தொடங்கினார்}.
\end{quote}

\begin{quote}
\emph{antat toḻilatipar oru putiya toḻiṟcālaiyait toṭaṅkiṉār.}
\end{quote}

\begin{quote}
That businessman started a new factory.
\end{quote}

\ta{தொழில்} \textquotedblleft{}a trade\textquotedblright{} + \ta{அதிபர்} \textquotedblleft{}a head, a chief\textquotedblright{}: the head of a trade.

\begin{cousinweb}
\textbf{\ta{தொழிலதிபர்}} (\emph{toḻilatipar}) joins the native \ta{தொழில்} (\emph{toḻil}), \textquotedblleft{}work, a trade\textquotedblright{} (DEDR 3524), and \ta{அதிபர்} (\emph{atipar}), \textquotedblleft{}a chief\textquotedblright{}, from Sanskrit \emph{adhipa}, \textquotedblleft{}a ruler\textquotedblright{} (Monier-Williams).
\end{cousinweb}

\begin{grammarlens}[title={What you've built}]
One more word for describing what you lived and felt.
\end{grammarlens}

\subsection*{Your turn}
\begin{itemize}
\raggedright
  \item \emph{Say it:} \emph{toḻilatipar}
  \item \emph{Say it:} \emph{toḻilatipar}, once more
  \item \emph{Say it:} \emph{antat toḻilatipar oru putiya toḻiṟcālaiyait toṭaṅkiṉār.} — That businessman started a new factory.
  \item \emph{Recall:} say the Tamil for imports, then the Tamil for business, then say \emph{toḻilatipar} again
\end{itemize}

\begin{tcolorbox}[breakable,colback=teal!4,colframe=teal!35!black,title={Before you move on}]
What does \ta{தொழிலதிபர்} mean? (\textquotedblleft{}A businessman; an industrialist\textquotedblright{}.) Say it once more.
\end{tcolorbox}
\section[\texorpdfstring{\ta{கூட்டாளி} (kūṭṭāḷi)}{கூட்டாளி (kūṭṭāḷi)}]{\ta{கூட்டாளி} (kūṭṭāḷi) — a partner; an associate}

\label{lesson:TA-C511-kuttali}



\begin{quote}
Before the new one: say the Tamil for business, then the Tamil for a businessman.
\end{quote}

\subsection*{You'll want to know: \ta{கூட்டாளி}}
\textbf{\ta{கூட்டாளி}} — \emph{kūṭṭāḷi} — \textquotedblleft{}a partner; an associate\textquotedblright{}.

\begin{quote}
\ta{நானும்} \ta{என்} \ta{கூட்டாளியும்} \ta{சேர்ந்து} \ta{ஒரு} \ta{கடை} \ta{நடத்துகிறோம்}.
\end{quote}

\begin{quote}
\emph{nāṉum eṉ kūṭṭāḷiyum cērntu oru kaṭai naṭattukiṟōm.}
\end{quote}

\begin{quote}
My partner and I run a shop together.
\end{quote}

Plural: \ta{கூட்டாளிகள்}, \emph{kūṭṭāḷikaḷ}.

\ta{கூட்டு} \textquotedblleft{}joining\textquotedblright{} + -\ta{ஆளி} \textquotedblleft{}one who …\textquotedblright{}. The plural is \ta{கூட்டாளிகள்}.

\begin{cousinweb}
\textbf{\ta{கூட்டாளி}} (\emph{kūṭṭāḷi}) is native Dravidian: \ta{கூட்டு} (\emph{kūṭṭu}), \textquotedblleft{}joining, partnership\textquotedblright{}, from \ta{கூடு} (\emph{kūṭu}), \textquotedblleft{}to come together\textquotedblright{} (DEDR 1882), with -\ta{ஆளி}, \textquotedblleft{}one who …\textquotedblright{}.
\end{cousinweb}

\begin{grammarlens}[title={What you've built}]
One more word for describing what you lived and felt.
\end{grammarlens}

\subsection*{Your turn}
\begin{itemize}
\raggedright
  \item \emph{Say it:} \emph{kūṭṭāḷi}
  \item \emph{Say it:} \emph{kūṭṭāḷi}, once more
  \item \emph{Say it:} \emph{nāṉum eṉ kūṭṭāḷiyum cērntu oru kaṭai naṭattukiṟōm.} — My partner and I run a shop together.
  \item \emph{Recall:} say the Tamil for business, then the Tamil for a businessman, then say \emph{kūṭṭāḷi} again
\end{itemize}

\begin{tcolorbox}[breakable,colback=teal!4,colframe=teal!35!black,title={Before you move on}]
What does \ta{கூட்டாளி} mean? (\textquotedblleft{}A partner; an associate\textquotedblright{}.) Say it once more.
\end{tcolorbox}
